% ============================================================================
%  SECCIÓN 3.18: USO DE WIDGETS DE REACT NATIVE Y DISEÑO MEDIANTE LAYOUTS
% ============================================================================

\section{Uso de widgets de React Native y dise\~no mediante layouts}

Para la construcci\'on de las interfaces de la aplicaci\'on \nombreApp{}, se
utilizaron los componentes nativos de React Native que equivalen a los widgets
de Flutter descritos en la gu\'ia del avance. A continuaci\'on se presenta un
an\'alisis de los widgets empleados y de los patrones de dise\~no de layouts
aplicados en las pantallas.

\subsection{Widgets de React Native utilizados}

Los widgets de React Native son componentes de interfaz que se renderizan
directamente en el dispositivo m\'ovil. En la siguiente tabla se presenta la
relaci\'on entre los widgets de Flutter recomendados en la gu\'ia del avance y
sus equivalentes en React Native, junto con los componentes adicionales que se
emplearon en el proyecto.

\begin{table}[H]
  \centering
  \caption{Relaci\'on de widgets Flutter vs React Native}
  \contenidoConFuente{%
    \begin{tablaAPA}
      \begin{tabular}{@{}p{3.4cm}p{4.0cm}p{8.2cm}@{}}
        \toprule
        \textbf{Widget Flutter} & \textbf{Equivalente React Native} &
        \textbf{Uso en el proyecto} \\\
        \midrule
        \texttt{TextField} & \texttt{TextInput} & Campos de formulario
        (email, contrase\~na, nombre) en las pantallas de login y registro. \\\
        \addlinespace
        \texttt{Text} & \texttt{Text} & T\'itulos, subt\'itulos, mensajes,
        etiquetas y textos informativos en todas las pantallas. \\\
        \addlinespace
        \texttt{ElevatedButton} & \texttt{TouchableOpacity} & Botones de
        acci\'on: \textit{``Comenzar''}, \textit{``Iniciar sesi\'on''},
        \textit{``Crear cuenta''}, \textit{``Cerrar sesi\'on''}. \\\
        \addlinespace
        \texttt{Column} & \texttt{View} con \texttt{flexDirection: 'column'} &
        Organizaci\'on vertical de elementos en todas las pantallas. \\\
        \addlinespace
        \texttt{Row} & \texttt{View} con \texttt{flexDirection: 'row'} &
        Disposici\'on horizontal: barra de navegaci\'on, footer de enlaces,
        bot\'on de logout con icono. \\\
        \addlinespace
        \texttt{Padding} & \texttt{padding*} en \texttt{StyleSheet} &
        Espaciado interno de contenedores, campos y botones. \\\
        \addlinespace
        \texttt{Container} & \texttt{View} con estilos de fondo y borde &
        Contenedores de avatar, campos de formulario y cajas de mensajes. \\\
        \addlinespace
        -- & \texttt{KeyboardAvoidingView} & Ajuste autom\'atico de la
        interfaz cuando aparece el teclado virtual. \\\
        \addlinespace
        -- & \texttt{ScrollView} & Desplazamiento del contenido cuando
        excede el tama\~no de la pantalla. \\\
        \addlinespace
        -- & \texttt{ActivityIndicator} & Indicador de carga durante el
        env\'io de formularios (login y registro). \\\
        \addlinespace
        -- & \texttt{Alert} & Di\'alogo de confirmaci\'on para el cierre
        de sesi\'on. \\\
        \bottomrule
      \end{tabular}
    \end{tablaAPA}%
  }{Elaboraci\'on propia en base al c\'odigo fuente de la aplicaci\'on.}
\end{table}

Todos los widgets se importan desde el m\'odulo \texttt{react-native}, que
proporciona componentes nativos optimizados para cada plataforma. Adem\'as, se
utiliz\'o la librer\'ia \texttt{@expo/vector-icons/Ionicons} para los
\'iconos de la barra de navegaci\'on inferior y de la pantalla de perfil.

\subsection{Uso de estilos con StyleSheet}

Los estilos de cada componente se definen mediante el objeto
\texttt{StyleSheet.create()} de React Native, que permite declarar estilos
tipados y optimizados para cada pantalla. A continuaci\'on se presenta un
ejemplo del estilo del bot\'on de la pantalla de bienvenida.

\begin{lstlisting}
const styles = StyleSheet.create({
  button: {
    backgroundColor: colors.white,
    paddingVertical: spacing.md + 2,
    borderRadius: borderRadius.full,
    alignItems: 'center',
    shadowColor: colors.black,
    shadowOffset: { width: 0, height: 4 },
    shadowOpacity: 0.15,
    shadowRadius: 8,
    elevation: 6,
  },
  buttonText: {
    color: colors.primary,
    fontSize: fontSize.lg,
    fontWeight: fontWeight.bold,
    letterSpacing: 0.5,
  },
});
\end{lstlisting}

\subsection{Dise\~no mediante layouts (flexbox)}

El sistema de layout de React Native se basa en \textbf{Flexbox}, un modelo de
dise\~no que permite organizar los elementos de forma flexible y responsiva. A
diferencia de Flutter, que utiliza \texttt{Row} y \texttt{Column} como
widgets separados, en React Native se controla la direcci\'on del layout a
trav\'es de la propiedad \texttt{flexDirection} del contenedor \texttt{View}.

\subsubsection{Layout vertical (equivalente a Column)}

El layout vertical es el patr\'on m\'as utilizado en las pantallas de la
aplicaci\'on. Se aplica mediante \texttt{flexDirection: 'column'} (valor por
defecto de \texttt{View}) y permite organizar los elementos de arriba hacia
abajo.

\begin{lstlisting}
// Layout vertical: encabezado, campos y botón
<View style={styles.container}>
  <View style={styles.header}>
    <Text style={styles.title}>Bienvenido</Text>
    <Text style={styles.subtitle}>Inicia sesión para continuar</Text>
  </View>
  <View style={styles.form}>
    <TextInput ... />
    <TextInput ... />
    <TouchableOpacity ... />
  </View>
</View>
\end{lstlisting}

Este patr\'on se aplica en la pantalla de login (encabezado + formulario +
footer), en la pantalla de registro (encabezado + 4 campos + bot\'on + enlace)
y en la pantalla de perfil (avatar + nombre + email + bot\'on de cerrar sesi\'on).

\subsubsection{Layout horizontal (equivalente a Row)}

El layout horizontal se utiliza para elementos que se disponen de izquierda a
derecha, como el footer de enlaces en las pantallas de login y registro, y la
barra de navegaci\'on inferior.

\begin{lstlisting}
// Layout horizontal: footer con enlace a registro
<View style={styles.footer}>
  <Text style={styles.footerText}>¿No tienes cuenta?</Text>
  <Link href="/auth/register" asChild>
    <TouchableOpacity>
      <Text style={styles.link}> Registrate</Text>
    </TouchableOpacity>
  </Link>
</View>
\end{lstlisting}

\begin{lstlisting}
// Layout horizontal: botón de cerrar sesión con icono
<TouchableOpacity style={styles.logoutButton} onPress={handleSignOut}>
  <Ionicons name="log-out-outline" size={20} color={colors.error} />
  <Text style={styles.logoutText}>Cerrar sesión</Text>
</TouchableOpacity>
\end{lstlisting}

\subsubsection{Uso de flex para distribuir espacio}

La propiedad \texttt{flex} permite distribuir el espacio disponible entre los
elementos hijos. En la pantalla de bienvenida se utiliza \texttt{flex: 1} en
el contenedor de contenido para que ocupe todo el espacio vertical disponible,
empujando el bot\'on \textit{``Comenzar''} hacia la parte inferior.

\begin{lstlisting}
content: {
  flex: 1,              // Ocupa todo el espacio disponible
  alignItems: 'center',  // Centra horizontalmente
  justifyContent: 'center', // Centra verticalmente
  zIndex: 1,
},
\end{lstlisting}

\subsubsection{Contenedores con Card (equivalente a Container)}

Los contenedores visuales tipo tarjeta se implementan con \texttt{View} +
estilos de fondo, borde y sombra. El ejemplo m\'as representativo es el avatar
de la pantalla de perfil, que utiliza un contenedor circular con sombra.

\begin{lstlisting}
logoWrapper: {
  width: 140,
  height: 140,
  borderRadius: borderRadius.full,
  backgroundColor: colors.white,
  alignItems: 'center',
  justifyContent: 'center',
  shadowColor: colors.black,
  shadowOffset: { width: 0, height: 8 },
  shadowOpacity: 0.2,
  shadowRadius: 16,
  elevation: 10,
},
\end{lstlisting}

\subsubsection{Listas y datos est\'aticos}

En esta etapa del desarrollo, los datos se presentan de forma est\'atica en
cada pantalla. La pantalla de perfil muestra la informaci\'on del usuario
obtenida del contexto de autenticaci\'on, y las pantallas de b\'usqueda y
favoritos presentan contenido placeholder. La integraci\'on con datos din\'amicos
desde Supabase se realizar\'a en avances posteriores.
